% ===================================================================
%  नकसा नं. १५ — बजार। --- खाए के चीज।
%
%  Traduction, faite en 2026, en bhojpouri d'aujourd'hui, sur l'ido
%  de texto/io. Regles de la colonne : voir 10-nakasa-01.tex. Une
%  seule numerotation. 28 blocs, 102 renvois — le plus dense du
%  livret —, quinze paires a retourner. L'ido coupe l'alinea 17 sur
%  un « suite » de feuillet ; la traduction n'herite pas de la coupe.
%
%  LA POMME DE TERRE (61) EST « आलू » ET NON « बटाटा », ET CELA
%  PRECISE TOUT CE QUE LA COLONNE A DIT DU PORTUGAIS. Elle a releve
%  cinq mots portugais en quatorze tableaux — मेज la table, पादरी le
%  cure, मिस्त्री le macon, अस्पताल l'hopital, आया la bonne d'enfants
%  — et les cinq figuraient dans la liste de la colonne gujaratie ;
%  le tableau 11 en avait conclu que la couche portugaise n'etait pas
%  cotiere mais indienne. Ce tableau-ci montre l'exception, et elle
%  est reguliere. La colonne gujaratie avait aussi releve « બટાટા »
%  la pomme de terre (de batata) et « કોબી » le chou (de couve). Ces
%  deux-la n'ont PAS traverse : la plaine dit आलू et गोभी. LA
%  DIFFERENCE EST NETTE ET ELLE TIENT EN UNE PHRASE : les mots
%  portugais qui ont franchi le sous-continent sont ceux qui
%  nommaient une chose SANS NOM — une table haute, un pretre
%  catholique, un maitre-macon, un hopital, une nourrice a gages ; ceux
%  qui nommaient un legume se sont arretes a la cote, parce que la
%  plaine avait un vieux mot a recycler. आलू nommait une igname en
%  sanscrit et a ete repris pour le tubercule neuf. ON N'EMPRUNTE PAS
%  CE QU'ON PEUT DIRE AVEC CE QU'ON A — troisieme verification de la
%  meme loi, apres पनडुब्बी au tableau 10 et डिब्बा au tableau 12.
%
%  LA VIANDE DE BŒUF (94) EST INTERDITE AU BIHAR DEPUIS 1955, ET LA
%  COLONNE L'ECRIT. C'est la troisieme fois que le releve rencontre
%  cette situation et la troisieme fois qu'il tranche de meme. La
%  colonne gujaratie l'avait rencontree A CE MEME TABLEAU 15 et au
%  meme renvoi (94) : le Gujarat punit l'abattage de la vache de la
%  reclusion a perpetuite depuis 2017. Le tableau 13 de cette
%  colonne-ci avait ecrit un Saint-Emilion que le Bihar interdit
%  depuis 2016. Le Bihar Preservation and Improvement of Animals Act
%  est de 1955 et il est plus vieux que les deux autres. LA LANGUE A
%  LES MOTS — गोमांस est dans la loi, dans les jugements et dans les
%  journaux ; c'est le faire qui est interdit, pas le nommer.
%
%  ET « पनीर » N'EST PAS DU FROMAGE. Le mot existe, il est ordinaire,
%  et il ne convient a aucun des quatre fromages de l'alinea 2 : le
%  paneer se caille a l'acide et non a la presure, ne s'affine pas, ne
%  moisit pas et ne fond pas. Un camembert (10) et une meule de
%  gruyere (11) sont une autre chose. La colonne ecrit donc « चीज »,
%  l'anglais. C'est le cas du pain au tableau 6 pris par l'autre
%  bout : la, la langue avait fabrique डबलरोटी pour marquer la
%  difference ; ici elle n'a rien fabrique, et il a fallu emprunter.
%  Le levantin, lui, n'avait pas eu ce probleme : جبنة couvre les deux.
%
%  Enfin l'epicerie (32) est « किराना के दुकान », et किराना est le mot
%  meme que la colonne gujaratie avait releve — « કરિયાણું », les
%  denrees seches, celles qui se transportent. Le mot a traverse
%  l'Inde ; le legume, non.
% ===================================================================

%%K t15-apar-1 apar
\VUsaut{4.0mm}
\VUtitre{15.0pt}{60}{नकसा नं. १५}
\VUsaut{3.0mm}

%%K t15-tit-1 sub
\VUpk{11.4pt}{40}{बजार। --- खाए के चीज।}
\VUsaut{3.0mm}

%%K t15-01-1 p
१. --- जे बेवसाई हमनी के खाए-पिए के जरूरी सामान देलें, ओहमें से जादेतर
\VUgras{ढँकल बजार} के \VUgras{बड़का छाजन} \textsuperscript{(1)} के नीचे भा
आसपास के गली में जुटल रहेलें।

%%K t15-02-1 p
२. --- \VUgras{सूअर के कसाई} \textsuperscript{(2)}, जे \VUgras{हैम}
\textsuperscript{(3)}, \VUgras{खून के सॉसेज} \textsuperscript{(4)},
\VUgras{सॉसेज} \textsuperscript{(5)}, \VUgras{ओझरी के सॉसेज}
\textsuperscript{(6)} आ \VUgras{सूअर के चरबी} \textsuperscript{(7)} बेचेलें,
\VUgras{मक्खन आ चीज बेचे वाली} \textsuperscript{(8)} के बगल में ठाढ़
बाड़ें, जेकर बसता \VUgras{हॉलैंड के चीज} \textsuperscript{(9)} से भरल बा
जेकर छिलका लाल बा, \VUgras{कामाम्बेर चीज} \textsuperscript{(10)} से,
\VUgras{ग्रुयेर के चक्का} \textsuperscript{(11)} से आ ताजा \VUgras{मक्खन
के पिंड} \textsuperscript{(12)} से।

%%K t15-03-1 p
३. --- ओकरा सामने एगो \VUgras{तरकारी बेचे वाली} \textsuperscript{(13)} आपन
ग्राहकिन लोग के सुंदर \VUgras{टमाटर} \textsuperscript{(14)},
\VUgras{फूलगोभी} \textsuperscript{(15)}, \VUgras{सलाद} \textsuperscript{(16)},
\VUgras{बंदागोभी} \textsuperscript{(17)}, \VUgras{कोहड़ा}
\textsuperscript{(19)}, \VUgras{खरबूजा} \textsuperscript{(18)} आ इहाँ तक कि
\VUgras{कुकुरमुत्ता} \textsuperscript{(20)} देखावत बाड़ी। बाकिर एगो
\VUgras{बीयर के छोकड़ा}, जे आपन \VUgras{पहुँचावे के गाड़ी}
\textsuperscript{(21)} — जवना पर \VUgras{बीयर के पीपा} \textsuperscript{(22)}
लदल बा — ठीक ओकर बसता के सामने रोक देले बा, ओकर ग्राहक लोग के लगे आवे
से रोकत बा। एगो मालिन ओकरा खातिर \VUgras{हाथीचक} \textsuperscript{(23)} के
एगो टोकरी ले आवत बाड़ी।

%%K t15-tit-2 sub
\VUsaut{4.0mm}
\VUpk{10.2pt}{40}{बेचल आ खरीदल।}
\VUsaut{3.0mm}

%%K t15-04-1 p
४. --- बजार में बहुते चहल-पहल बा, काहेकि \VUgras{बोली}
\textsuperscript{(24)} के बेरा आ गइल बा। \VUgras{घरनी लोग}
\textsuperscript{(26)}, जे इहाँ सौदा करे आवेली, चलत-चलत \VUgras{ओझरी बेचे
वाली} \textsuperscript{(25)} के दुकान के सामने रुक जाली, \VUgras{ओझरी} भा
\VUgras{बछड़ा के कपार} खरीदे खातिर।

%%K t15-05-1 p
५. --- एगो तखिन \VUgras{बावर्ची} \textsuperscript{(27)} \VUgras{मछरी बेचे
वाली} \textsuperscript{(28)} से एगो बहुते सुंदर \VUgras{टूना} के भाव-ताव
करत बाड़ें, आ ऊ आपन मरजी से भाव बढ़ावत बाड़ी। ओकरा लगे \VUgras{बाम, सोल
मछरी, ट्राउट} आ \VUgras{कार्प} के बड़का खेप आइल बा। ओकर \VUgras{कॉड} आ
\VUgras{सीप} बहुते ताजा बा।

%%K t15-tit-3 sub
\VUsaut{4.0mm}
\VUpk{10.2pt}{40}{सस्ता आ महँग माल।}
\VUsaut{3.0mm}

%%K t15-06-1 p
६. --- \VUgras{मुरगी बेचे वाली} \textsuperscript{(29)}, जे ओकरा बगल में
बइठल बाड़ी, बहुते ईमानदार बाड़ी। जब ऊ \VUgras{पीरू, हंस, बतख} भा सादा
\VUgras{मुरगी} बेचेली, त खाली सही दाम माँगेली।

%%K t15-07-1 p
७. --- चउक के कोना पर, पहिला मंजिल पर, कुछ दिन से एगो \VUgras{कपड़ा
बेचे वाला} \textsuperscript{(30)} बइठल बा, जेकरा लगे खरीदार मेहरारू लोग
के हरमेसा भरोसा रहेला कि \VUgras{मेजपोस, नैपकिन, एप्रन} आ
\VUgras{पोंछन} के बहुते सुंदर तरह-तरह मिल जाई। ओही मंजिल पर एगो
\VUgras{छूरी बनावे वाला} \textsuperscript{(31)} आपन \VUgras{कारखाना} खोलले
बा। ऊ छाजन के बेचे वाली लोग के \VUgras{छूरी} तेज करेला आ मालिन लोग
खातिर सबसे कड़ा \VUgras{फौलाद} के \VUgras{हँसुआ} बनावेला।

%%K t15-tit-4 sub
\VUsaut{4.0mm}
\VUpk{10.2pt}{40}{किराना के दुकान।}
\VUsaut{3.0mm}

%%K t15-08-1 p
८. --- ओकर कारखाना के नीचे, कुछ दिन से, एगो बड़का \VUgras{खाना के
दुकान} खुलल बा। ऊ आउर ऊ सादा आ मामूली \VUgras{किराना के दुकान}
\textsuperscript{(32)} नइखे रह गइल जवन पहिले रहे, जे खाली रोज के काम के
चीज बेचत रहे — \VUgras{चाय} \textsuperscript{(33)}, \VUgras{चॉकलेट}
\textsuperscript{(34)}, \VUgras{चीनी के चक्का} \textsuperscript{(37)} आ
\VUgras{साबुन} \textsuperscript{(38)}। इहाँ हर चीज बिकेला —
\VUgras{सूखल बिस्कुट}, \VUgras{डिब्बाबंद} \textsuperscript{(35)},
\VUgras{मीठ सराब} \textsuperscript{(36)}, \VUgras{रम, कोंयाक, किर्श,
सौंफ के सराब} आ \VUgras{कुराकाओ}। इहाँ \VUgras{सिकार} भी मिलेला :
\VUgras{खरहा} \textsuperscript{(39)}, \VUgras{थ्रश} \textsuperscript{(40)},
\VUgras{तीतर} \textsuperscript{(41)}, \VUgras{बटेर} \textsuperscript{(42)},
\VUgras{जंगली बतख} \textsuperscript{(43)} भा \VUgras{फीजेंट}
\textsuperscript{(44)}, आ ओतने आसानी से परदेसी फर भी, जइसे \VUgras{केरा}
\textsuperscript{(46)} आ \VUgras{अनन्नास}।

%%K t15-09-1 p
९. --- किराना के \VUgras{छोकड़ा} \textsuperscript{(45)}, जे बहुते मददगार
बा, जवन माँगीं ऊ देवे के तइयार बा : \VUgras{करौंदा} भा \VUgras{बिही} के
\VUgras{मुरब्बा} \textsuperscript{(47)}, \VUgras{चरबी} \textsuperscript{(48)},
\VUgras{खीरा के अचार} \textsuperscript{(49)} भा नून लागल \VUgras{सारडीन}
\textsuperscript{(50)}।

%%K t15-09-2 p
ई \VUgras{ग्राहकिन} \textsuperscript{(51)} लोग खातिर बहुते सुभीता बा, जे
सबसे मामूली चीज के बगल में सबसे दुर्लभ \VUgras{नया फसल} आ सबसे सवाद
वाला \VUgras{फर} पा जाली : रसदार \VUgras{नासपाती} \textsuperscript{(52)},
\VUgras{आलूबोखारा} \textsuperscript{(53)}, \VUgras{अंगूर} \textsuperscript{(54)},
\VUgras{आड़ू} \textsuperscript{(55)}, \VUgras{बादाम} \textsuperscript{(56)}
आ \VUgras{अखरोट} \textsuperscript{(57)}।

%%K t15-tit-5 sub
\VUsaut{4.0mm}
\VUpk{10.2pt}{40}{ग्राहक लोग।}
\VUsaut{3.0mm}

%%K t15-10-1 p
१०. --- \VUgras{चाउर} \textsuperscript{(58)} आ \VUgras{मसूर}
\textsuperscript{(59)} \VUgras{धुँआइल हेरिंग} \textsuperscript{(60)} के पेटी
के बगल में बा ; \VUgras{आलू} \textsuperscript{(61)} आ \VUgras{सेम}
\textsuperscript{(63)} \VUgras{सेब} \textsuperscript{(62)}, \VUgras{अंजीर}
\textsuperscript{(64)}, \VUgras{सूखल आलूबोखारा} \textsuperscript{(65)} आ
\VUgras{हेजलनट} \textsuperscript{(66)} के लगे बा। एह दुकान में ताजा
\VUgras{अंडा} \textsuperscript{(67)} आ \VUgras{जैतून} \textsuperscript{(68)}
भी मिलेला ; एहीसे \VUgras{टोकरी} \textsuperscript{(70)} आ \VUgras{सौदा के
जाली} \textsuperscript{(73)} पलक झपकते भर जाला।

%%K t15-11-1 p
११. --- एह बढ़िया दुकान के \VUgras{कॉफी के दाना} \textsuperscript{(72)}
के \VUgras{नौकरानी} \textsuperscript{(69)} आ रसोइनिन लोग के बीच सचहूँ नाँव
हो गइल बा, काहेकि बुढ़उ \VUgras{किराना वाला} \textsuperscript{(71)} खुदे
ओकरा के बहुते जतन से भूँजेलें।

%%K t15-tit-6 sub
\VUsaut{4.0mm}
\VUpk{10.2pt}{40}{नजारा।}
\VUsaut{3.0mm}

%%K t15-12-1 p
१२. --- सब लोग बजार में सौदा करे ना आवे। जवन गली छाजन तक ले जाले ओकर
रंगीन आवाजाही में सबसे अलग-अलग तरह के लोग लउकेला। एगो भल
\VUgras{बूचड़खाना के छोकड़ा} \textsuperscript{(75)} के बगल में, जे आपन
सामने \VUgras{ठेला} \textsuperscript{(74)} ठेलत बा, ई रहलें दू गो छुट्टी पर
आइल सिपाही, जे आपन चमकत वरदी देखावे में बहुते फूलल बाड़ें :
\VUgras{हुसार} \textsuperscript{(76)}, नीला \VUgras{दोलमान} आ पंख वाला
\VUgras{साको} में, आ \VUgras{जुआव के नायक}, फुलल जाँघिया में आ
\VUgras{तुर्की टोपी} पहिनले।

%%K t15-13-1 p
१३. --- \VUgras{नानबाई} \textsuperscript{(80)}, जे \VUgras{नानबाई के
दुकान} \textsuperscript{(78)} के देहरी पर ठाढ़ बाड़ें, अभिए आपन
\VUgras{डबलरोटी} \textsuperscript{(79)} के आटा गूँथले बाड़ें आ भट्ठी से ऊ
\VUgras{क्रोसां} \textsuperscript{(81)}, \VUgras{पाव} \textsuperscript{(82)}
आ \VUgras{बड़का रोटी} \textsuperscript{(83)} निकालले बाड़ें जवन सो-केस में
बा।

%%K t15-14-1 p
१४. --- नानबाई के दुकान के ऊपर एगो कुसल \VUgras{बियाढ़न के सिलाई करे
वाली} \textsuperscript{(84)} रहेली, जे कई गो \VUgras{कढ़ाई करे वाली}
\textsuperscript{(85)} के काम पर रखले बाड़ी। ओकरा बगल में एगो
\VUgras{धोबिन आ इस्तिरी करे वाली} \textsuperscript{(86)} रहेली, जे
\VUgras{बियाढ़न} \textsuperscript{(88)} के धो के आ सुखा के ओकरा में कलफ
लगावेली, आ \VUgras{इस्तिरी} \textsuperscript{(87)} से ओकरा के इस्तिरी
करेली।

%%K t15-tit-7 sub
\VUsaut{4.0mm}
\VUpk{10.2pt}{40}{मांस, मछरी आ तरकारी।}
\VUsaut{3.0mm}

%%K t15-15-1 p
१५. --- \VUgras{कसाई के दुकान} \textsuperscript{(89)} में, जवन
भुइयाँ-तल पर बा, हर तरह के \VUgras{मांस} बिकेला — \VUgras{भेड़ के मांस}
\textsuperscript{(90)}, \VUgras{गोमांस} \textsuperscript{(94)} भा
\VUgras{बछड़ा के मांस} \textsuperscript{(92)} — \VUgras{जाँघ, बीफस्टेक,
भूँजल मांस} आ \VUgras{पसुली} के रूप में। \VUgras{कसाई के छोकड़ा}
\textsuperscript{(93)} ई सब मांस काटेला आ \VUgras{गूदा वाला हाड़}
\textsuperscript{(96)} अलग धरेला। कसाई के दुकान के केवाड़ पर एगो
\VUgras{सीप बेचे वाली} \textsuperscript{(97)} बाड़ी जे \VUgras{सीप}
\textsuperscript{(98)} बेचेली। केहू चाहे त ऊ ओकरा के खोल भी देली।

%%K t15-16-1 p
१६. --- ई सब बेवसाई रखसा भा जगहा के भाड़ा देलें ; एहीसे \VUgras{पुलिस
वाला} \textsuperscript{(100)} ओह बेचारा \VUgras{ठेला पर तरकारी बेचे वाली}
\textsuperscript{(99)} लोग के बेदरदी से खदेड़ेला, जे आपन ठेला पर
\VUgras{गाजर, मूली, सलगम, गंधना} भा \VUgras{सतावर के मुट्ठा, हरियर मटर,
पालक, चूका, हलीम} आ \VUgras{अजमोद} ढो के ओह लोग के मुकाबिला करेली।

%%K t15-tit-8 sub
\VUsaut{4.0mm}
\VUpk{10.2pt}{40}{पुलिस से बचीं !}
\VUsaut{3.0mm}

%%K t15-17-1 p
१७. --- जे लइका \VUgras{लहसुन} \textsuperscript{(101)} आ \VUgras{पियाज}
बेचेला ऊ नीक से जानेला कि ओकरो के बजार लगे आपन माल बेचे के हक नइखे। ऊ
दउड़ के भागेला, एह खतरा के साथे कि ऊ \VUgras{बड़का झींगा} आ
\VUgras{लॉबस्टर} बेचे वाला \VUgras{बेचे वाला} \textsuperscript{(102)} के
\VUgras{केकड़ा, नदी के झींगा} आ \VUgras{झींगा} के टोकरी गिरा दी।

%%K t15-18-1 p
१८. --- सब लोग हिलत-डोलत बा आ भयानक सोर करत बा ; बाकिर एहसे ओह
\VUgras{सीसा वाला} \textsuperscript{(103)} के आवाज साफ सुने में कवनो बाधा
ना पड़े जे गली में फेरी लगावत बा, ना ओह \VUgras{कोयला वाला}
\textsuperscript{(104)} के हाँक, जे अभिए कुछ पल खातिर आपन ठेला छोड़ के
एगो \VUgras{कोयला के बोरा} \textsuperscript{(105)} पहुँचावे गइल बा।
